\begin{afterword}
\summaryhead

The post answers the homework set three months earlier, to trace the thinking that produces debates about free will. People who hear that physics is deterministic picture ``Me'' and ``Physics'' as rival causes of the future: if physics fixes the future, there is no room left for me. The post redraws the diagram with ``Me'' inside ``Physics.'' The mistake is hard to see, it says, because no one can picture atoms making up a person the way one sees fingers making up a hand.

The author then names a position, ``Requiredism'': choice, control and moral responsibility, sensibly understood, require determinism, at least in patches, because no one can plan and act amid utter chaos. If the future were not determined by physics, it would be determined by something, and that something would include you. Physics includes our decisions and does not explain them away; the confusion comes from our cognitive algorithms.

\responsehead

The post's diagrams make a clear point. Asking whether I or physics decides the future is like asking whether a hand or its fingers picks something up. The post's reply to the argument that physics would make reasoning mere ``quarks bopping around'' is the one Elizabeth Anscombe made to C.~S. Lewis: good reasons remain good reasons, whatever causal story we tell about the person who has them.

The position is not new, and it is not a third option beside compatibilism. If free will requires determinism, it is compatible with it, so ``Requiredism'' is a strong form of compatibilism. Hume argued in 1748 that actions not caused by a person's character reflect neither credit nor blame on that person, and Hobart's 1934 paper put the thesis in its title. The post does not claim novelty. In a 2009 survey, most philosophers accepted or leaned toward compatibilism.

The case against the rivals is thinner than the case for the position. The post puts the alternative to determinism as ``utter chaos.'' Libertarians in the literature do not hold that. Event-causal views such as Robert Kane's require only that some decisions be caused by the agent's reasons without being fixed by them. Agent-causal views say the agent determines the decision. The post's weakened formula, that the future must be ``determined by something,'' is one an agent-causal libertarian also accepts. What excludes them is the clause ``some law, some order,'' which is asserted. The objection the post points toward, that undetermined choices would be a matter of luck, is one these views are often thought vulnerable to. The post does not take up their replies.

The post's explanation of the confusion, an implicit causal network in the brain, is its answer to the homework; like the network in ``Neural Categories,'' it is a guess about design, not evidence about brains.

In short: a clear picture of an old compatibilist answer, under a new name, defended against ``utter chaos,'' which is not what the libertarian views in the literature hold.
\end{afterword}
